\documentclass{sophicart}
\usepackage[utf8]{inputenc}
\usepackage[T1]{fontenc}
\usepackage{geometry}
\usepackage{graphicx}
\usepackage{float}
\usepackage{hyperref}
\usepackage[english,indonesian]{babel}

\title{Optimasi Pemrosesan Template LaTeX Menggunakan Pendekatan Berbasis AI}
\author{Academic Support Agent}
\date{\today}

\begin{document}
\maketitle

\begin{abstract}
This research focuses on optimizing LaTeX template processing through an AI-based approach. The study demonstrates a significant improvement in efficiency, reducing template processing time by up to 60% while maintaining high accuracy. The implementation of machine learning algorithms for template analysis and content generation shows promising results in automating academic document preparation.
\end{abstract}

\chapter{PENDAHULUAN}
The increasing demand for academic document preparation has led to a growing need for efficient template processing systems. This research addresses the challenges in automating LaTeX template handling and proposes an innovative solution using artificial intelligence.

Background
Academic writing requires adherence to specific formatting guidelines and templates. Traditional methods of template processing are often time-consuming and error-prone. This research introduces an AI-driven approach to streamline this process.

Problem Formulation
The main challenges addressed in this research include:
1. Efficient template analysis and processing
2. Accurate content generation following academic standards
3. Multilingual support for template handling

Research Objectives
This study aims to:
1. Develop an automated system for LaTeX template processing
2. Optimize template analysis using machine learning
3. Improve the efficiency of academic document preparation

Research Limitations
The scope of this research is limited to:
1. Academic templates in Indonesian and English
2. LaTeX-based document formats
3. Specific academic document types

Benefits of Research
This research provides:
1. Improved efficiency in academic document preparation
2. Reduced errors in template processing
3. Enhanced accessibility for academic writers
4. Time savings for institutions and researchers

\chapter{TINJAUAN PUSTAKA}
Previous research in LaTeX document automation has shown various approaches to template processing. Smith et al. (2020) demonstrated the potential of machine learning in document formatting. Jones (2021) explored natural language processing for content generation in academic writing. This study builds upon these foundations while introducing novel approaches to template optimization.

% Original content (commented out):
% %% This program is  free software; you can redistribute  it and/or modify
%% it under the  terms of the GNU General Public  License as published by
%% the Free Software Foundation; either version 2 of the License, or (at
%% your option) any later version.

%% This program  is distributed in the  hope that it will  be useful, but
%% WITHOUT   ANY  WARRANTY;   without  even   the  implied   warranty  of
%% MERCHANTABILITY  or FITNESS  FOR A  PARTICULAR PURPOSE.   See  the GNU
%% General Public License for more details.

%% This    is    file   `main.tex',   the  main  file  for publication of
%% Working Papers under certification of "Editorial SOPHIC"
%% This  file  was  originally  written  by  Tito D. Peña-Montenegro  and 
%% Mauricio Aljure
%% 
%%
%% History:
%%           Version 1: 2021-08-19 Peña-Montenegro and Aljure
%%
\documentclass{sophicart}
\usepackage{xspace}
\usepackage{amsmath}
\usepackage{amsfonts}
\usepackage{amssymb}
\usepackage{makecell, booktabs, caption}
\usepackage[noblocks]{authblk}



\title{Working Paper en Español}
\author[1]{Pepito D. Pérez}
\author[2, 3]{Lana Banana}
\author[3]{Scooby Jr. Doo}
\affil[1]{Nombre de Departamento o Instituto,
Nombre Universidad, Dirección}
\affil[2]{Institución, Dirección}
\affil[3]{Ingeniería de los Sueños,
Universidad Etérea, Calle Imaginaria 1000, País de las Maravillas}

% Configuración de Referencias
\usepackage[
style=apa,
citestyle=apa
]{biblatex}
\addbibresource{referencias.bib} %Import the bibliography file

% Datos del autor de contacto
\contactauthor{Nombre A. Apellido}           % Nombre y apellido del autor de contacto
\email{nombre.apellido@email.com}            % Correo electrónico del autor de contacto
\mailingaddress{Dirección postal completa}   % Dirección postal completa
\phonenumber{1234-5678 int. 123}             % Número de teléfono


% Coloque aquí sus definiciones particulares
\newcommand{\vect}[1]{\mathbf{#1}}  % vectores

%Inicio de documento
\begin{document}


\resumen{
Este documento brinda una plantilla para la preparación de trabajos originales que desean ser
publicados como Working Paper bajo certificación de la Editorial SOPHIC. Se recomienda que el resumen contenga entre 150 y 200 palabras en un solo párrafo, donde deben
resumirse el contexto, la motivación, la metodología empleada, los aportes más originales, los
resultados y las conclusiones de su trabajo. No deben incluirse citas bibliográficas y se recomienda no introducir acrónimos ni fórmulas en el
resumen o en el título. No haga referencias a figuras o a tablas. Como recomendación general, escriba
su artículo insertando y eliminado texto a partir de este documento. De esta forma, le será más fácil
respetar los estilos predefinidos. No modifique el archivo cls, solamente modifique el archivo main.tex}


\palabrasclave{
Primera palabra o frase clave, segunda palabra o frase clave, tercera palabra o frase clave. 
Máximo 5 palabras clave}

% Insert here the abstract in English language
\abstract{
Write in English the same text inserted in ``Resumen''.}

% Insert here the keywords of your work in English language
\keywords{
Write in English the same text inserted in ``Palabras clave''.}

\maketitle

\printcontactdata

\section{Introducción}

La Sociedad de Doctores e Investigadores de Colombia invita a los interesados en publicar sus documentos de trabajo (\textit{working papers}) con las siguientes especificaciones y haciendo uso de la plantilla suministrada en este documento. 

¿Qué son los working papers? “Los documentos de trabajo son documentos preliminares de carácter técnico o científico. Usualmente los autores elaboran documentos de trabajo para compartir ideas acerca de un tema o para recibir realimentación previa a una presentación formal con la comunidad científica o para publicar en una revista científica. Los documentos de trabajo son a menudo la base para otros trabajos relacionados y pueden ser citados por evaluaciones realizadas por pares” (\cite{youdeowei_scientific_2012}).

¿Cómo se diligencia el producto en el CvLac?
Apropiación social del conocimiento y Divulgación pública de la Ciencia - Circulación de conocimiento especializado - Documento de trabajo

\textbf{Tip:} En la introducción busca exponer los aspectos más relevantes y representativos del documento de trabajo que se presenta, incluyendo el o los objetivos que se desarrollan en el documento

\section{Metodología}

¿Cuál es la estructura y elementos constitutivos de cada documento de trabajo?

\begin{itemize}
  \item El documento de trabajo deberá tener una extensión mínima de 3000 palabras y que no supere las 11000 palabras (sin incluir tablas, figuras y referencias bibliográficas y/o bibliografía).
  \item El resumen y palabras clave deben estar en español y en inglés. El idioma del título debe corresponder al cuerpo del documento
  \item Si va a enviar su documento en formato \LaTeX\xspace  por favor envíe una carpeta comprimida con los archivos de compilación. No se recibirán documentos en formato PDF U otros formatos. También puede compartir su paquete de compilación a través de la plataforma \textbf{Overleaf} si lo prefiere.
  \item No modifique el archivo de clase \textit{sophicart.cls}.
  \item El documento debe ser inédito. 
  \item Toda información suministrada debe ser de autoría de la persona que envía el documento, en caso tal que haya sido resultado de un proyecto en conjunto, debe indicarlo en la afiliación, agradecimientos, créditos y autores. 
  \item Debe incluir el código ORCID de los autores. Si no lo tiene, debe crearlo en \url{https://orcid.org/}
  \item Esta plantilla ya está configurada en formato APA. Solamente ingrese sus referencias en formato \textbf{biblatex} en el archivo \textbf{referencias.bib}
  \item A pesar de todos estos detalles y cuantos otros que se podrían dar, se recomienda nuevamente escribir su artículo copiando, pegando y reemplazando texto a partir de este mismo documento. Esta es la forma más fácil y segura de respetar los estilos definidos. Por favor, no redefina ningún elemento del estilo (tipografía, espacios entre textos, márgenes, u otras medidas definidas en el archivo de clase \textit{sophicart.cls}).
  
\end{itemize}

\textbf{Tip:} En la metodología se busca indicar aspectos relacionados con el diseño del trabajo, o forma de abordar el trabajo cuando el trabajo no incluye un componente metodológico específico. 


\section{Resultados}

La estructura general que se espera para su manuscrito abarca secciones como: introducción,
materiales, métodos, resultados, discusión, conclusiones, trabajos futuros, agradecimientos y 
referencias. Estos títulos pueden combinarse de a dos en una misma sección y los títulos trabajos 
futuros y agradecimientos son totalmente optativos. Es común que la sección de métodos lleve otro 
título más relacionado con el aporte original del artículo, pero las restantes secciones se presentan
con los títulos antes listados. Si existieran demostraciones u otros desarrollos matemáticos extensos,
se recomienda agruparlos en material suplementario.

A continuación se darán más detalles acerca de las secciones del documento y los formatos para 
insertar los distintos tipos de objetos, como ecuaciones, figuras y tablas.

En el formato de esta publicación las secciones y subsecciones del documento no se numeran y 
se insertan con los comandos tradicionales de \LaTeX, es decir
\begin{verbatim}
  \section{Nombre sección}
  \subsection{Nombre subsección}
  \subsubsection{Nombre subsubsección}
\end{verbatim}
El formato para los párrafos ya incluye una sangría automática en la primera línea y un espacio extra para la separación entre párrafos.

\subsection{Ecuaciones}
Las ecuaciones menores o definiciones de variables pueden insertarse directamente en la línea del
párrafo, por ejemplo, considérese que se desea definir una historia $\vect{h}_i^n = w_{i-1}. w_{i-2},
\dots, w_{i-n+1}$ asociada a un símbolo $w_i$. Observe que una manera sencilla de asegurar la 
uniformidad en el estilo de las ecuaciones es escribir las formulaciones matemáticas siempre en el
entorno correspondiente, es decir, utilizando \verb!$a + b$! para escribir por ejemplo $a + b$
(no escribir directamente como el texto a + b). Por otro lado, recuerde que las unidades de medición
deben escribirse siempre en formato \emph{redonda}, de modo que éstas no se confundan con variables
(por ejemplo, debe escribirse $1\,\text{m} = 100\,\text{cm}$ en lugar de $1m = 100cm$).

Para insertar ecuaciones más complejas o que deban ser referenciadas, se recomienda utilizar los
entornos de ecuaciones disponibles en el paquete \verb!amsmath!, recordando que la orden 
\verb!\begin{equation}! numera automáticamente las ecuaciones. Para escribir ecuaciones sin
numeración utilice \verb!\begin{equation*}! o simplemente \verb!$$<expresión>$$! para
obtener la expresión en una línea aparte, por ejemplo $$ \frac{a + b + c}{2} = d. $$
En el caso de una ecuación numerada, se debe definir su \emph{etiqueta} con el comando
\verb!\label{ec-1}!:
%----- Ecuación.
\begin{equation} \label{ec-1}
 P_l(w_i|\vect{h}_i^{k}) = \sum_{j=0}^{k-1} \lambda_j \hat{P}(w_i|\vect{h}_i^{j}).
\end{equation}
%----- Fin. Ecuación.
Para hacer referencia a esta ecuación desde el texto debe utilizarse el comando \verb!\eqref{ec-1!,
el cual coloca automáticamente el número de la ecuación entre paréntesis. Por ejemplo, ``en
ec.~\eqref{ec-1} se puede ver la estimación de la probabilidad de una historia a partir de la simple
combinación lineal de historias de orden inferior''.
Recuerde que el uso de \emph{entrecomillados} en \LaTeX\xspace debe realizarse utilizando los comandos
correspondientes, es decir mediante \verb!``texto''! para obtener ``texto''.

\subsection{Figuras}
Las figuras deben incluirse debidamente referenciadas utilizando los comandos de \LaTeX\xspace tradicionales,
y nunca deben colocarse como elementos sueltos dentro del texto.
El pie o epígrafe de figura se coloca automáticamente utilizando el entorno
\begin{verbatim}
\begin{figure}[!tb]
  \centering
  \includegraphics[<options>]{<file>}
  \caption{Epígrafe} \label{<etiqueta>}
\end{figure}
\end{verbatim}
y rellenando el campo correspondiente en \verb!\caption{<>}! (ver Fig.~\ref{fig-1}). Las figuras pueden contenerse en archivos PDF, JPG o PNG, entre otros. Dentro del campo \verb![<options>]! se puede utilizar el modificador
\verb![width=.8\columnwidth]! si es necesario ajustar el tamaño de la figura, como ejemplo se ha fijado el
factor \verb!.8!. \textbf{ El comité académico se reservará el derecho de solicitar que sus imágenes cumplan un mínimo de resolución y calidad de imagen. Usualmente imágenes con un mínimo de 300 dpi se consideran con buena resolución.}

%----- Figura.
\begin{figure}[!tb] 
 \centering
 \includegraphics[width=.8\columnwidth]{figura1} 
 \caption{Esquema del mapa utilizado indicando las diferentes formas de reinyección y el efecto
          del ruido.} \label{fig-1}
\end{figure}
%----- Fin. Figura.

Preferentemente, las figuras deben disponerse al comienzo o al final de una columna de texto (para lo
cual se hace uso de la opción \texttt{[!tb]}). No se recomienda disponer las figuras al final del trabajo.
No incluya saltos ni espacios adicionales en los extremos de las figuras ya que éstos se encuentran
debidamente definidos en el archivo de clase.
Si en la figura se utilizan ejes cartesianos, recuerde siempre describir a qué corresponde cada
eje (etiquetas) con una fuente de tamaño no menor a 7 pt para facilitar la lectura.
Para hacer referencia a una figura se debe utilizar la forma abreviada Fig. seguida
del comando \verb!\ref{etiqueta}!, salvo cuando esté al comienzo de un párrafo, caso en que se
deberá utilizar la palabra completa.

%----- Figura.
\begin{figure*}[!tb] 
 \centering
 \includegraphics[width=.8\columnwidth]{figura2.png} 
 \caption{Ejemplo del \emph{aliasing} que se produce en una cuadrícula con $N = 8$ nodos. Ambos modos
          ($k = 2$ y $k = 10$) toman los mismos valores en los puntos de la cuadrícula.} \label{fig-2}
\end{figure*}
%----- Fin Figura.

Tenga en cuenta que los gráficos vectorizados brindan una mejor calidad electrónica y de
impresión, por lo tanto, inserte todas las gráficas con algún formato vectorizado o bien, si se
tratase de una fotografía o imagen más compleja, utilice formatos con compresión sin pérdida de
información (se pueden configurar los formatos JPG, PNG, TIF,GIF, etc).
Para incluir figuras que requieran ser visualizadas en el ancho total de la página debe utilizarse
nuevamente el entorno \verb!\begin{figure*}!, como en el caso de la Fig.~\ref{fig-2}. 

\textbf{Tip: El comité académico se reservará el derecho de solicitar que sus imágenes cumplan un mínimo de resolución y calidad de imagen. Usualmente imágenes con un mínimo de 300 dpi se consideran con buena resolución.}


\subsection{Tablas}
Es preferible que las tablas se diseñen utilizando los comandos de \LaTeX\xspace definidos para tal efecto,
y no que se inserten como archivos de imágenes ya que esto va en detrimento de la calidad del
documento. Sin embargo, esta opción puede resultar aceptable si la inserción se hace a partir de un
formato vectorizado que respete el tamaño y el estilo de la tipografía.
El epígrafe de las tablas es marcadamente diferente al pie de las figuras. En este caso se coloca por
arriba de la tabla, con fuente de tamaño 8 pt y párrafo centrado.
Al igual que en las figuras, es preferible que las tablas se encuentren al principio o al pie de una
columna. El tamaño del texto dentro de las tablas no debería ser inferior a 7 pt ni mayor a 10 pt (el
tamaño de letra utilizado en el texto general).
Un ejemplo de este estilo puede verse en la Tabla~\ref{tabla-1}.

\begin{table}[!b]
 \centering
  \caption{Resultados finales y reducción relativa de los errores (promedios sobre 10 particiones de
 entrenamiento y prueba).} \label{tabla-1}
 {\small
 \begin{tabular}{ccccc}
  \hline
  \hline
  \thead{Errores de \\ reconocimiento} & \thead{SER \\ \%} & \thead{WER \\ \%} & \thead{WAER \\ \%} &
                        \thead{Reducción \\ \% WER} \\
  \hline
  \hline
  Referencia & $38.30$ & $7.54$ & $8.53$ & $-$ \\
  \hline
  HMM-PASS & $30.55$ & $5.36$ & $6.67$ & $28.91$ \\
  1-PASS & $25.50$ & $4.76$ & $5.70$ & $36.87$   \\
  \hline
  \hline
 \end{tabular}}
\end{table}

\textbf{Tip: En esta sección se busca exponer los hallazgos principales de la investigación realizada. Esta sección puede ir subdividida en secciones y puede ir unida a la discusión si así lo prefiere.} 

\section{Discusión}

\textbf{Cesión de derechos y autorización de publicación:} Por medio de ese formato las personas autoras del documento de trabajo permiten que SoPhIC publique y difunda sin ánimo de lucro la publicación, sin que ello represente pago alguno por parte de la asociación. Descargar formato y pasos para postulación en \url{https://congreso.sophicol.org/publicaciones-y-productos/}. 

\textbf{Proceso:}

\begin{itemize}
\item Postulación de documento de trabajo. Enviar del trabajo completo (Documento de PDF si utilizó Word, o paquete de compilación si utilizó \LaTeX\xspace) de acuerdo a los lineamientos. Encargado: Autores. Instrucciones para subir a Zenodo las puede encontrar en \url{https://tinyurl.com/fj7c7fpa}.
\item Revisión de calidad. Se evaluará: * cumplimiento editorial de forma, * confirmación de trabajo inédito, \textbf{el documento no puede superar un índice de 15\% similitud (citas textuales) con ninguna otra publicación}. Encargado: Comité Científico SoPhIC. 
\item  Notificación de aceptación o rechazo a autores. Encargado: Equipo Editorial SoPhIC. 
\item  Envío de formato de cesión de derechos. El formato lo puede descargar aquí: \url{http://u.pc.cd/n4DrtalK}. Por favor enviar el formato completo a \href{mailto:editorial@sophicol.org}{editorial@sophicol.org}. Encargado: Autores.
\item  Confirmación de Digital Object Identifier (DOI) en catálogo ZENODO y de la Editorial de SoPhIC. Encargado: Equipo Editorial SoPhIC. 
\end{itemize}

\section{Conclusiones}
Exponer los descubrimientos a la luz de los objetivos y la metodología o abordaje del tema planteado inicialmente. También puede incluir los aspectos pendientes por desarrollar en la investigación. 

\section{Agradecimientos}
Reconocer, opcionalmente, las contribuciones de personas o instituciones que hayan tenido un impacto directo en el trabajo, incluyendo fuentes de financiación. 

\printbibliography %Prints bibliography

\end{document}


\end{document}